\begin{figure}[ht]
\centering
\begin{tikzpicture}[x=1cm, y=1cm, font=\small]

% Sample size n required to reach 80% power at alpha = 0.05 for a
% two-sample t-test, plotted against standardised effect size
% d = delta/sigma. Approx n_per_group = 16 / d^2.

% axes
\draw[->, thick] (0, 0) -- (8.5, 0)
  node[below, font=\scriptsize] {standardised effect size $d = \delta / \sigma$};
\draw[->, thick] (0, 0) -- (0, 5.5)
  node[left, font=\scriptsize, rotate=90, yshift=0.6cm, anchor=south]
  {$n$ per group};

% x-axis ticks for d = 0.1, 0.2, ..., 0.8
\foreach \d/\xc in {0.1/1, 0.2/2, 0.4/4, 0.6/6, 0.8/8} {
  \draw (\xc, 0) -- (\xc, -0.1)
    node[below, font=\scriptsize] {\d};
}
% y-axis ticks: 10, 100, 1000 (log scale)
\foreach \nval/\yc in {10/1, 100/3, 1000/5} {
  \draw (0, \yc) -- (-0.1, \yc)
    node[left, font=\scriptsize] {\nval};
}

% n = 16 / d^2 on log10 y axis: y = log10(16/d^2) - 1, mapped to plot
% units (log10(10) -> y=1, log10(1000) -> y=5, so y = 2*(log10(n)-1)
% + 1 means y plot = log10(n)*2 - 1)
% We sample d at fine increments and plot.
\draw[thick, engblue, smooth]
  plot[domain=0.10:0.85, samples=80]
  (\x*10, {2*(ln(16/(\x*\x))/ln(10)) - 1});

% Annotate three key points
\node[circle, fill=engred, inner sep=1.5pt, label={right:{$d=0.2,\ n\!\approx\!400$}}]
  at (2, {2*(ln(16/(0.2*0.2))/ln(10)) - 1}) {};
\node[circle, fill=engred, inner sep=1.5pt, label={right:{$d=0.5,\ n\!\approx\!64$}}]
  at (5, {2*(ln(16/(0.5*0.5))/ln(10)) - 1}) {};
\node[circle, fill=engred, inner sep=1.5pt, label={right:{$d=0.8,\ n\!\approx\!25$}}]
  at (8, {2*(ln(16/(0.8*0.8))/ln(10)) - 1}) {};

\end{tikzpicture}
\caption[Sample size scales as the inverse square of effect size]{For
a two-sample $t$-test at $\alpha = 0.05$ and power $0.80$, the
required sample size per group is approximately $n \approx 16/d^{2}$
where $d = \delta/\sigma$ is the standardised effect size (Cohen's
$d$). Halving the effect size quadruples the required sample. The
plot shows the rule-of-thumb curve on a logarithmic $n$ axis; an
effect of $d = 0.2$ (small) demands roughly $400$ per group, $d =
0.5$ (medium) about $64$, $d = 0.8$ (large) about $25$. Studies
chasing small effects with small samples are the statistical
mechanism behind much of the failure-to-replicate pattern of
section 8.8.}
\label{fig:vol01:ch08:effect-size-vs-n}
\end{figure}
